% Generated draft: requires mathematical and source review.
\documentclass[11pt]{article}
\usepackage[margin=25mm]{geometry}
\usepackage{fontspec}
\setmainfont{Latin Modern Roman}
\usepackage{amsmath,amssymb,mathtools,unicode-math}
\setmathfont{Latin Modern Math}
\usepackage{enumitem,needspace,xurl,hyperref}
\setlength{\parindent}{0pt}
\setlength{\parskip}{5pt}
\setlength{\emergencystretch}{3em}
\newcommand{\N}{\mathbb{N}}
\newcommand{\Z}{\mathbb{Z}}
\newcommand{\Q}{\mathbb{Q}}
\newcommand{\R}{\mathbb{R}}
\newcommand{\C}{\mathbb{C}}
\newcommand{\F}{\mathbb{F}}
\newcommand{\A}{\mathbb{A}}
\newcommand{\PP}{\mathbb{P}}
\newcommand{\E}{\mathbb{E}}
\newcommand{\eps}{\varepsilon}
\newcommand{\dd}{\,\mathrm{d}}
\begin{document}
\section*{UnsolvedMath ID: 30000559}
\textbf{Source problem code:} catalogue entry
\subsection{Definitions and mathematical statement}
\textbf{Entropy-spectrum problem and arithmetic question.} Let \(\alpha:\Z^k\to\operatorname{Diff}^{1+\varepsilon}(M)\), \(\varepsilon>0\), act on a \((k+1)\)-manifold, with an ergodic invariant measure, no proportional Lyapunov exponents, and at least one positive-entropy element. What values can the entropies of the elements take? The supplied arithmetic conjecture asks whether these values are algebraic integers of degree at most \(k+1\). This arithmetic wording is retained literally, with its normalization issue explicit below.\par In a usual probability-measure interpretation, let \(\mu\) be invariant under the whole action and ergodic for that action, and put \(\mathcal E(\alpha,\mu)=\{h_\mu(\alpha(a)):a\in\Z^k\}\), with measure-theoretic entropy in a stated logarithm convention. This does not require every individual element to be ergodic. The positive-entropy hypothesis is \(\exists a_0\in\Z^k\) with \(h_\mu(\alpha(a_0))>0\). Under the necessary derivative integrability hypotheses, joint Lyapunov exponents are linear functionals on \(\R^k\); the stated nonproportionality condition compares these functionals. It must not silently be replaced by the stronger general-position condition on every collection of Lyapunov hyperplanes. The original refers to maximal-rank hypotheses of Theorem 1, which is not reproduced; the permitted \(k\), multiplicities, zero exponents, compactness, and measure assumptions require source checking.\par Literally, the arithmetic claim is that for every \(a\), the real number \(h_\mu(\alpha(a))\) is a root of a monic polynomial in \(\Z[T]\) with algebraic degree at most \(k+1\). Ordinary entropy is logarithmic and changes scale with the base of the logarithm. An assertion about algebraic integers, one about logarithms of algebraic numbers, and one about \(\exp(h_\mu(\alpha(a)))\) are different assertions. The source gives no normalization or explanation identifying them. In particular, the arithmetic wording must not be silently repaired by exponentiating the entropy or inserting a logarithm. This is a substantive verification need, not merely a typography issue. The record also does not explicitly distinguish measure entropy from topological entropy.\par The supplied triage checked 2026-08-21 labels the item open and says no authoritative resolution was located. Its later verification field reports recovering Theorem 1's hypotheses, Problem 2, and the following arithmetic restriction from the official report, but none of that surrounding text is included here. Both the restored wording and these source limitations are retained without an independent mathematical or current-status assessment.
\subsection{Short English statement}
Describe the entropy values of the specified maximal-rank action and verify the proposed arithmetic restriction. The literal algebraic-integer claim must be distinguished from a claim about logarithms or exponentials of entropy. The missing theorem also prevents replacing nonproportionality by stronger Lyapunov general-position assumptions.
\subsection{Sources}
\begin{itemize}\item Local UnsolvedMath registry, permanent ID 30000559, code catalogue entry; \emph{Entropy Spectra of Maximal-Rank Abelian Actions}, Oberwolfach Reports Open Problems; source context identifies Theorem 1 and Problem 2, neither reproduced in full.\item \emph{Geometric Group Theory, Hyperbolic Dynamics and Symplectic Geometry (2007)}, Oberwolfach Reports, DOI 10.4171/owr/2006/33, pp. 1991--2058. \url{https://doi.org/10.4171/owr/2006/33}. Proposed year and year tag: 2006; no proposer named.\item The same full report is the sole supplied literature source. No separate authoritative entropy-arithmetic theorem or normalization reference is given.\item Literature triage checked 2026-08-21, status \texttt{open}; record updated 2026-09-16. The \texttt{corrected\_verified} field reports restoring the theorem hypotheses and arithmetic restriction, not independently checked here.\end{itemize}
\end{document}
